
\noindent\textbf{命题~\ref{prop:7-2-275} 的证明收束。}
令
\[
\begin{aligned}
\widehat{\mathfrak C}'_1
={}&\widehat{\mathfrak C}_a(L'_1,\lambda'_1)
\vee_r\widehat{\mathfrak C}_Y(L'_1)
\vee_r\widehat{\mathfrak C}_d(L'_1,\lambda'_1),\\
\mathfrak H
={}&(\mathfrak H_a\circ
\mathfrak f^{\{\psi_{0,s}\}_s})
\vee_r(\mathfrak J\circ\mathfrak H_X)
\vee_r(\mathfrak H_d\circ\mathfrak f^\phi).
\end{aligned}
\]
这里
\[
\begin{aligned}
\mathfrak f^{\{\psi_{0,s}\}_s}\colon\;&
C(L_0,\lambda_0)
\longrightarrow C([0,1]_s\times L'_0,
(\lambda_0,\lambda'_0)),\\
\mathfrak f^\phi\colon\;&
C(L_0,\lambda_0)
\longrightarrow C(L_1,(\lambda_0,\lambda_1)).
\end{aligned}
\]
再令
\[
\begin{aligned}
\mathfrak C_1'^{(s=s_0)}
&=\widehat{\mathfrak C}'_{1,s=s_0},\\
\mathfrak C'_0
&=C([0,1]_s\times L'_0,\lambda'_0),\\
\mathfrak C'_1
&=C([0,1]_s\times L'_1,\lambda'_1),\\
\widehat{\mathfrak C}'_{1,r=r_0}
&=\mathfrak C'_1,\\
\mathbb J'_0
&=j_{[0,1]_s}\circ
\mathfrak f^{\operatorname{id}\times\phi'}.
\end{aligned}
\]
其中 $r_0=0,1$\translatornote{原书在端面等式中
把下标写为通用的 $i$；此处只构造
$\widehat{\mathfrak C}'_1$，
故按命题条件 (7.2.276.2)
限定为 $i=1$。另：原书本段先写无撇的
$\mathbb J_0$、随即使用带撇的
$\mathbb J'_0$，且与引理 7.2.342 的
$\mathbb J$ 字形相近；译文统一按本处实际使用的
$\mathbb J'_0$ 排写，见疑误 A-0746。}。

首先检验 $\mathfrak H$ 定义良好。
由 (7.2.311.2) 得
\[
\begin{aligned}
\operatorname{Eval}_{r=1}\circ\mathfrak H_a
\circ\mathfrak f^{\{\psi_{0,s}\}_s}
={}&j_{[0,1]_s}\circ\mathfrak E_a
\circ\mathfrak f_{\mathrm{ext}}^{\operatorname{id}\times\phi'}
\circ\mathfrak f^{\{\psi_{0,s}\}_s}.
\end{aligned}
\]
另一方面，图式~\ref{diag:7-2-28}、
\ref{diag:7-2-27} 的交换性给出
\[
\begin{aligned}
\operatorname{Eval}_{r=0}\circ\mathfrak J\circ\mathfrak H_X
&=j_{[0,1]_s}\circ\mathfrak E_a
\circ\operatorname{Eval}_{r=0}\circ\mathfrak H_X\\
&=j_{[0,1]_s}\circ\mathfrak E_a
\circ\mathfrak f_{\mathrm{ext}}^{\operatorname{id}\times\phi'}
\circ\mathfrak f^{\{\psi_{0,s}\}_s}.
\end{aligned}
\]
同样由图式~\ref{diag:7-2-28}、
\ref{diag:7-2-27} 的交换性，
\[
\begin{aligned}
\operatorname{Eval}_{r=1}\circ\mathfrak J\circ\mathfrak H_X
&=\mathfrak e_d\circ\mathfrak E_d
\circ\operatorname{Eval}_{r=1}\circ\mathfrak H_X\\
&=\mathfrak e_d\circ\mathfrak E_d
\circ\mathfrak f_{\mathrm{ext}}^{\{\psi_{1,s}\}_s}
\circ\mathfrak f^\phi.
\end{aligned}
\]
而 (7.2.337.1) 给出
\[
\operatorname{Eval}_{r=0}\circ\mathfrak H_d
\circ\mathfrak f^\phi
=\mathfrak e_d\circ\mathfrak E_d
\circ\mathfrak f_{\mathrm{ext}}^{\{\psi_{1,s}\}_s}
\circ\mathfrak f^\phi.
\]
故 $\mathfrak H$ 确实定义良好。

下面逐项核验 (7.2.276.1)--(7.2.276.6)。
前三项容易验证。对于第四项，计算得
\[
\begin{aligned}
\operatorname{Eval}_{r=0}\circ\mathfrak H
&=\operatorname{Eval}_{r=0}\circ\mathfrak H_a
\circ\mathfrak f^{\{\psi_{0,s}\}_s}\\
&=j_{[0,1]_s}\circ
\mathfrak f^{\operatorname{id}\times\phi'}
\circ\mathfrak e_a
\circ\mathfrak f^{\{\psi_{0,s}\}_s}\\
&=\mathbb J'_0\circ
\mathfrak g^{\{\psi_{0,s}\}_s}.
\end{aligned}
\]
第一、三步依定义，第二步由 (7.2.311.1) 得到。
于是 (7.2.276.4) 成立。
再计算
\[
\begin{aligned}
\operatorname{Eval}_{r=1}\circ\mathfrak H
&=\operatorname{Eval}_{r=1}\circ\mathfrak H_d
\circ\mathfrak f^\phi\\
&=\mathfrak e_d\circ
\mathfrak f^{\{\psi_{1,s}\}_s}
\circ j\circ\mathfrak f^\phi\\
&=\mathfrak g^{\{\psi_{1,s}\}_s}
\circ\mathfrak g^\phi.
\end{aligned}
\]
其中第一、三步依定义，第二步由 (7.2.337.2) 得到。
于是 (7.2.276.5) 成立。

最后计算
\[
\begin{aligned}
\operatorname{Eval}_{s=s_0}\circ\mathfrak H
={}&(\operatorname{Eval}_{s=s_0}\circ\mathfrak H_a
\circ\mathfrak f^{\{\psi_{0,s}\}_s})
\vee_r(\operatorname{Eval}_{s=s_0}\circ\mathfrak J
\circ\mathfrak H_X)\\
&\vee_r(\operatorname{Eval}_{s=s_0}\circ\mathfrak H_d
\circ\mathfrak f^\phi)\\
={}&(\mathcal F_a^{s=s_0}\circ
\mathfrak f^{\psi_{0,s_0}})
\vee_r\Bigl(\mathfrak J^{(s=s_0)}\circ
\bigl(\mathcal F_b^{s=s_0}
\vee_r\mathfrak f^{\{\Phi_{r,s_0}\}_r}
\vee_r\mathcal F_c^{s=s_0}\bigr)\Bigr)\\
&\vee_r(\mathcal F_d^{s=s_0}\circ\mathfrak f^\phi)\\
={}&\mathfrak g_+^{\{\Phi_{r,s_0}\}_r}.
\end{aligned}
\]
第一步依定义；第二步由 (7.2.311.3)、
(7.2.273)、(7.2.343.1)、(7.2.341)、
(7.2.340.2) 及 (7.2.337.3) 得到；
第三步亦依定义。
事实上，$\mathfrak g_+^{\{\Phi_{r,s_0}\}_r}$
是 (7.2.298.1)、(7.2.298.2)、
(7.2.280.2)\textup{（译者注：源书误写为
(7.2.280.3)；后者是右侧两段的总复合。）}、
(7.2.330.1)、(7.2.330.2) 这五段同伦的复合；
它们分别对应上式第三行胶合的五个因子。
由此得到 (7.2.276.6)，
命题得证。\qed
